\chapter{Overview}
\label{chap:introduction}

\section{About AirfoilTools}

\textbf{AirfoilTools} is a Python library of tools dedicated
to the 2D aerodynamic design and analysis of airfoils, propellers and other profiled sections.
It relies on Bézier curves of arbitrary degree and
multi-segment splines for geometric modeling, and integrates a coupling with the \emph{XFoil} aerodynamic solver
for performance analysis.

The application combines:

\begin{itemize}
    \item a \textbf{geometric engine} based on Bézier
        curves of arbitrary degree and multi-segment splines
        (BezierSpline) with configurable continuities (C\textsuperscript{0},
        C\textsuperscript{1}, C\textsuperscript{2});
    \item a \textbf{coupling with XFoil}, the reference
        aerodynamic solver for viscous subsonic 2D
        flows;
    \item an \textbf{interactive graphical interface} allowing
        airfoil editing, simulation control and
        comparative visualization of results;
    \item a \textbf{tracing mode}: loading a background image
        (photo, scan, drawing) to reconstruct an airfoil
        by overlaying the control points onto the silhouette, all
        of which can be saved in a \textbf{project} file
        (\chemin{.aftproj});
    \item a \textbf{NACA airfoil generator} (4 and 5 digits)
        directly from the interface;
    \item a \textbf{flap module}: generation of a profile
        deflected about a hinge axis, which can be simulated and compared
        to the original airfoil;
    \item \textbf{geometric comparison tools}: deviation
        between two airfoils (vertical or normal) and curvature
        plotting.
\end{itemize}

\screenshot{01_main_window.png}{1.0}{%
    General view of the application at startup. The current airfoil
    (NACA~2412, blue, converted to Bézier) is overlaid on the reference
    airfoil (NACA~2412, red); the airfoil with deflected flap
    (green) as well as the curvature and the deviation are displayed by
    default.}

\section{Who this manual is for}

This manual is intended for the \textbf{users} of the graphical
application. It covers all the features accessible
from the interface, from loading an airfoil to
the comparative analysis of aerodynamic polars.

Developers interested in the Python API of the library
will find the technical documentation in the source code
(Sphinx-format docstrings) and in the generated documentation
(\chemin{docs/} folder).

\begin{noteinfo}
No prior knowledge of theoretical aerodynamics
is required to use the application. However, a
general understanding of aerodynamic quantities (lift
coefficient $C_L$, drag coefficient $C_D$, lift-to-drag ratio $C_L/C_D$,
Reynolds number, angle of attack $\alpha$) makes
interpreting the results easier.
\end{noteinfo}

\section{Functional architecture}

The application is structured around \textbf{three tabs}
corresponding to as many logical steps:

\begin{enumerate}
    \item \menu{Airfoils} --- loading, editing and comparison of
        geometric airfoils. Two airfoils are handled simultaneously:
        an editable \emph{current airfoil} (in blue), and a
        \emph{reference airfoil} (in red) that serves as a
        comparison.
    \item \menu{XFoil Settings} --- entry of the simulation
        parameters: Reynolds range, angle-of-attack sweep,
        boundary layer parameters, panel options.
    \item \menu{Results} --- graphical presentation of the results
        as a \emph{configurable grid of analysis
        cells}, each displaying an independent plot
        (polar $C_L(\alpha)$, $C_D(\alpha)$, lift-to-drag ratio, pressure
        distribution $-C_p(x)$, boundary layer profile, etc.).
\end{enumerate}

\section{Airfoil representation}

AirfoilTools offers two airfoil representation modes:

\subsection*{Discrete points mode}

The airfoil is defined by a sequence of points in the
\textbf{Selig} convention (running from the trailing edge to the upper surface, then
the leading edge, then the lower surface, and back to the trailing edge).
This is the most common format of airfoil libraries
(for example \chemin{airfoiltools.com}).

\subsection*{Spline mode (multi-segment Bézier)}

The airfoil is defined by two \textbf{BezierSpline} (upper surface and
lower surface), each composed of several Bézier segments
joined with C\textsuperscript{1} continuity. This mode allows:

\begin{itemize}
    \item \emph{interactive editing} by moving the
        control points;
    \item \emph{on-demand resampling} with a freely
        chosen point density;
    \item the \emph{exact computation of the curvature and tangents}
        at any point;
    \item the \emph{splitting of a segment} into two at an arbitrary point
        (De Casteljau algorithm, mathematically exact).
\end{itemize}

Switching from one mode to the other is done via the
\menu{Edit $\rightarrow$ Convert to Spline} menu or by loading
a file in \chemin{.bspl} or \chemin{.bez} format.

\begin{astuce}
To fully benefit from the interactive editing
features and curvature plotting, it is recommended
to \textbf{convert your airfoils to Spline mode} as soon as they are
loaded. The following chapters detail this procedure.
\end{astuce}
